\section{Introduction}

Enterprise large language model (LLM) applications are often explored first as demonstrations: a system prompt, a set of retrieved documents, and a user interface that shows the intended interaction. We call a prototype prompt-dominant when important product behavior is encoded mainly as natural-language instructions or broad retrieval context, with code, data contracts, and validation artifacts still underdeveloped. Prompt-dominant development is useful for exploration and is related to vibe coding, a conversational AI-assisted software workflow popularized by Karpathy \citep{karpathy2025vibecoding,meske2025vibecoding}. Productization introduces a different requirement: each visible claim should be traceable to bounded sources, routed to the correct entity, regenerated under the same assumptions, and inspected through explicit artifacts. These concerns connect to prior work on hallucination and software-engineering controls for artificial intelligence (AI) systems \citep{ji2023hallucination,martinezfernandez2022seai,kreuzberger2023mlops}.

The study examines this transition through the reconstruction of an LLM-based investment briefing agent. The motivating prototype was developed during a ten-week AI leadership program hosted by CEO AI \citep{ceoai2026leadership}; the authors report the work under the AI Leadership Research Center affiliation associated with that program context. During the program, the authors built and demonstrated over ten exploratory AI product prototypes for executive- and client-facing scenarios, using Replit-hosted demonstrations to share interaction patterns with participants \citep{replit2026platform}. The investment-briefing prototype supplied useful interaction patterns, including mobile-first briefing cards, source links, voice-oriented flows, and follow-up questions. The prototype serves as motivating context for a reconstructed system with a different source, trace, and validation structure.

We propose a harness-engineering approach for converting a demonstration prototype into a traceable research system. In software engineering, a test harness refers to scaffolding code used to exercise lower-level code before the higher-level code that will ultimately invoke it is available \citep{isoiecieee2017vocabulary}. We extend the test-harness idea to enterprise LLM agents: the harness is the code-owned control layer around the model, responsible for source gates, routing rules, claim eligibility, answer contracts, trace logging, and validation. In the reconstructed architecture, manifests define what sources may be used, source-backed claims define what statements may enter runtime context, routing metadata binds questions to entities, answer contracts define the visible response, and traces record how the answer was assembled. Maintained wiki pages, following the LLM Wiki pattern, provide compact context while source manifests and source-backed claims remain authoritative \citep{karpathy2026llmwiki}.

The paper contributes a harness-engineering method for reconstructing prompt-dominant enterprise LLM prototypes into traceable LLM-agent architectures; a source-to-claim knowledge pipeline that separates raw documents, evidence records, runtime-eligible claims, maintained wiki context, and user-facing answers; a replaceable composition boundary that separates deterministic harness control from LLM phrasing; and a system-level validation design for checking traceability, routing, prompt-to-code migration, internal-output leakage prevention, runtime-interface behavior, and latency. We instantiate the method on a bounded public-data slice of five Korean corporate groups, 25 listed companies, and 113 source-backed runtime claims.
